\section{Worked certificates and existence of admissible runs}\label{art:hypothesis-examples}

This appendix exhibits explicit band certificates and full-BFGS trajectories for
Definition~\ref{art:band-certificate} and establishes general sufficient
conditions---together with a sharp counterexample---for the existence of
infinite admissible runs from prescribed initial data. Throughout, norms are
Euclidean.

\subsection{Explicit band certificates and BFGS trajectories}

\begin{proposition}[Degenerating trajectory and band certificate on $u^2+|v|$]\label{art:examples-mixed}
Let $F(u,v)=u^2+|v|$ and $K_b=\{x\in\R^2:F(x)\le b\}$ ($b>0$). Then $F$ is
finite and convex with unique minimizer $0$ and compact sublevel $K_b$, and for
every $0<\eta\le b$, Definition~\ref{art:band-certificate} holds on $K_b$ with
\[
 W_\eta(x)=xx^T,\qquad m_\eta=\eta^2,\quad C_\eta=2b,
 \quad L_\eta=0,\quad\delta_\eta=1.
\]
Moreover, for $b=2/3$ and $(c_1,c_2)=(1/4,1/2)$, the sequence
\[
 x_k=\bigl(0,\tfrac23(-\tfrac12)^k\bigr)^T,\qquad
 g_k=\bigl(0,(-1)^k\bigr)^T,\qquad
 H_k=\begin{pmatrix}1&0\\0&2^{-k}\end{pmatrix},\qquad \alpha_k=1
\]
is an infinite admissible run in $K_{2/3}$ starting from $H_0=I$ and converging
to $(0,0)$, along which $\lambda_{\min}(H_k)=2^{-k}\to0$.
\end{proposition}
\begin{proof}
At any differentiable point $x=(u,v)^T\in K_b$ with $F(x)\ge\eta$,
$g=(2u,\operatorname{sign}v)^T$ gives $g^TW_\eta(x)g=(2u^2+|v|)^2\ge F(x)^2\ge\eta^2$.
For any differentiable $x=(u,v)^T$ and $z=(r,w)^T$ in $K_b$ with $t=y^T(z-x)>0$,
the two-case Cauchy--Schwarz calculation in the proof of
Corollary~\ref{art:mixed-benchmark} gives $(y^Tz)^2\le2F(z)t\le2bt$, and
$z=x+s$ gives $\Pi_sz=\Pi_sx$, verifying~\eqref{art:visibility}--\eqref{art:transport}.
Along the displayed sequence, $x_k$ is differentiable with $g_k\ne0$,
$s_k=-H_kg_k=(0,-(-1/2)^k)^T$, $y_k=-2g_k$, $d_k=2^{-k}$,
$F(x_k)-F(x_{k+1})=\tfrac13d_k\ge\tfrac14d_k$, $t_k=2d_k$, and~\eqref{art:update}
maps $\operatorname{diag}(1,2^{-k})$ to $\operatorname{diag}(1,2^{-k-1})$.\qed
\end{proof}

\begin{proposition}[Full-rank band certificate in $\R^n$]\label{art:examples-full-rank}
For $F(x)=\tfrac12\norm{x}^2$ on $\R^n$ ($n\ge1$) and $b>0$, $K_b=\{F\le b\}$ is
compact with unique minimizer $0$, and for every $0<\eta\le b$,
Definition~\ref{art:band-certificate} holds with the rank-$n$ constant field
\[
 W_\eta(x)=I,\qquad m_\eta=2\eta,\quad C_\eta=1,
 \quad L_\eta=0,\quad\delta_\eta=1.
\]
For every $x_0\ne0$, $x_k=2^{-k}x_0$, $g_k=x_k$, $H_k=I$, $\alpha_k=1/2$ is an
infinite admissible run in $K_{F(x_0)}$ for $(c_1,c_2)=(1/4,1/2)$ converging to
$0$.
\end{proposition}
\begin{proof}
Here $g^TIg=\norm{x}^2=2F(x)\ge2\eta$, $y=s$ gives $y^TIy=t$, and constant
$W_\eta=I$ satisfies~\eqref{art:transport} with $L_\eta=0$. Along the stated
history, $y_k=s_k=-\tfrac12x_k$ leaves $H_k=I$ invariant under~\eqref{art:update},
$F(x_k)-F(x_{k+1})=\tfrac34d_k\ge\tfrac14d_k$, and $g_{k+1}^Ts_k=\tfrac12g_k^Ts_k$.\qed
\end{proof}

\subsection{A general existence criterion and curved-interface examples}

\begin{lemma}[Existence of differentiable Armijo--weak-Wolfe steps and runs]\label{art:examples-run-existence}
Let $F:\R^n\to\R$ be finite, convex, and bounded below, and suppose that for
every differentiability point $x$ of $F$ and every direction $p\ne0$, the set
\[
 \{t\in\R:F\text{ is not differentiable at }x+tp\}
\]
is countable (where differentiability refers to $F$ on $\R^n$). Then at any
differentiable $x$ with $g=\nabla F(x)\ne0$, any $H\succ0$, and any
$0<c_1<c_2<1$, the search ray $x+tp$ ($p=-Hg$, $t>0$) contains a differentiable
point $x+\alpha p$ satisfying both~\eqref{art:aw} with strict inequalities. If,
in addition, $\nabla F(x)\ne0$ at every differentiability point of $F$, then an
infinite admissible run exists from every differentiable $x_0$, every
$H_0\succ0$, and every $0<c_1<c_2<1$.
\end{lemma}
\begin{proof}
Set $\phi(t)=F(x+tp)$, $D=g^Tp<0$, and $\sigma(t)=(\phi(t)-\phi(0))/t$ for
$t>0$. Differentiability at $x$ gives $\lim_{t\downarrow0}\sigma(t)=D$, while
boundedness below gives $\liminf_{t\to\infty}\sigma(t)\ge0>c_1D$. Since
$D<c_2D<c_1D<0$, continuity of $\sigma$ yields a nonempty open interval of
$t>0$ on which $c_2D<\sigma(t)<c_1D$. Picking $\alpha$ in that interval outside
the countable nondifferentiable set gives strict Armijo $\phi(\alpha)<\phi(0)+c_1\alpha D$,
and convexity of $\phi$ gives $\nabla F(x+\alpha p)^Tp=\phi'(\alpha)\ge\sigma(\alpha)>c_2D$,
hence strict weak Wolfe. By Lemma~\ref{art:run-invariants}, $H_+\succ0$; if no
differentiable point is stationary, repeating this choice inductively yields an
infinite admissible run.\qed
\end{proof}

\begin{proposition}[Curved two-quadratic family with global convergence and run existence]\label{art:examples-two-quadratics}
For any $a,b,c>0$, define on $\R^2$
\[
 F_{a,b,c}(u,v)=\max\{au^2+bv^2+v,\;cu^2+bv^2-v\},
\]
whose switching locus is the parabola $v=\tfrac12(c-a)u^2$ (a line only when
$a=c$). Then $F_{a,b,c}$ has unique minimizer $(0,0)$, compact sublevels,
satisfies Corollary~\ref{art:planar-convergence} on every positive sublevel and
Corollary~\ref{art:planar-max} at $(0,0)$, admits an infinite admissible run
from every differentiable $x_0$, every $H_0\succ0$, and every $0<c_1<c_2<1$, and
every such run converges in value and iterates to $(0,0)$.
\end{proposition}
\begin{proof}
Both branches have Hessian at least $2\min\{a,b,c\}I\succ0$ and average to
$\tfrac12(a+c)u^2+bv^2$, giving unique minimizer $(0,0)$, compact sublevels, and
strictly convex positive contours. Along the tie curve $v=\tfrac12(c-a)u^2$, the
two branch gradients $(2au,2bv+1)^T$ and $(2cu,2bv-1)^T$ have determinant
$-2u(a+c+b(c-a)^2u^2)$, which vanishes only at $(0,0)$; for $u\ne0$ they are
linearly independent and their affine hull misses $0$. Thus
$\mathcal C_{\rm aff}=\{(0,0)\}$ and Corollary~\ref{art:planar-convergence}
applies on every positive sublevel, while at $(0,0)$ the active gradients
$(0,1)^T$ and $(0,-1)^T$ are affinely independent with weights $1/2$. Finally,
$(0,0)$ is nondifferentiable (so every differentiable point is nonstationary),
and along any line through a differentiable point the branch difference
$(a-c)u^2+2v$ is a nonzero polynomial of degree at most two with at most two
roots, so Lemma~\ref{art:examples-run-existence} guarantees run existence from
every differentiable start.\qed
\end{proof}

\begin{corollary}[Run existence on the three-quadratic nonradial example]\label{art:examples-nonradial-runs}
For the objective $F$ of Theorem~\ref{art:nonradial-example}, every
differentiable $x_0$ with $F(x_0)<1/100$ (for instance $x_0=(1/1000,0)^T$, where
$F(x_0)=1001/10^6<1/100$), every $H_0\succ0$, and every $0<c_1<c_2<1$ admit an
infinite admissible run, and every such run converges to $(0,0)$.
\end{corollary}
\begin{proof}
Here $F(u,v)\ge u^2+v^2$ has unique minimizer $(0,0)$, where $F$ is
nondifferentiable, so every differentiable point has $\nabla F(x)\ne0$. For each
pair $i<j$, the quadratic difference $f_i-f_j$ has positive-definite or
indefinite nonzero quadratic and linear parts such that its restriction to any
nonconstant line is a nonzero polynomial of degree at most two: indeed,
$f_0-f_1=u-v+\tfrac12u^2-v^2$, $f_0-f_2=2u+v-\tfrac12u^2-\tfrac12v^2$, and
$f_1-f_2=u+2v-u^2+\tfrac12v^2$, none of which can vanish identically along a
line (for $f_0-f_2$, the quadratic part $-\tfrac12(u^2+v^2)$ never vanishes on a
nonzero direction; for $f_0-f_1$ and $f_1-f_2$, any null direction $d_1=\pm\sqrt2\,d_2$
or $d_2=\pm\sqrt2\,d_1$ of the quadratic part has irrational slope ratio and
therefore nonzero pairing with the integer linear part). Every line therefore
meets the pairwise tie sets in finitely many points, and
Lemma~\ref{art:examples-run-existence} together with
Theorem~\ref{art:nonradial-example} completes the proof.\qed
\end{proof}

\subsection{Why theorems quantify over admissible runs: an obstructed initial ray}

\begin{proposition}[A coercive homogeneous objective with no differentiable first step]\label{art:examples-no-step}
The function $G(u,v)=\max\{u,-u+v,-u-v\}$ is finite, convex, coercive, and
positively homogeneous of degree one with unique minimizer $(0,0)$ (hence covered
by Theorem~\ref{art:homogeneous}). Nevertheless, at any $x_0=(\varepsilon,0)^T$
($\varepsilon>0$), $G$ is differentiable, yet with $H_0=I$ the search ray
$x_0-\alpha H_0\nabla G(x_0)$ ($\alpha>0$) contains \emph{no} differentiable
point satisfying weak Wolfe for any $c_2<1$.
\end{proposition}
\begin{proof}
Since $G(u,v)\ge\max\{u,-u\}=|u|$ and $2G(u,v)\ge u+(-u\pm v)=\pm v$ gives
$2G(u,v)\ge|v|$, $G$ is coercive with unique minimizer $(0,0)$. At
$x_0=(\varepsilon,0)^T$, $\nabla G(x_0)=(1,0)^T$, so the search ray is
$(\varepsilon-\alpha,0)^T$ ($\alpha>0$). For $0<\alpha<\varepsilon$,
$\nabla G(\varepsilon-\alpha,0)=(1,0)^T$, so weak Wolfe would require
$-1\ge-c_2$, impossible for $c_2<1$; for $\alpha\ge\varepsilon$, the two branches $-u+v$ and $-u-v$ both evaluate to
$\alpha-\varepsilon\ge\varepsilon-\alpha$ with distinct gradients $(-1,1)^T$ and
$(-1,-1)^T$, so $(\varepsilon-\alpha,0)^T$ is nondifferentiable.\qed
\end{proof}

